\documentclass[../../main.tex]{subfiles}
% \documentclass[a4paper,10pt]{ltjsarticle}
\begin{document}
\rhead{京都大学後期$1998$年$3$番}
\setcounter{equation}{0}
\begin{oframed}
  \documentclass[../../main.tex]{subfiles}
\begin{document}

$\mathrm{A}_1$，$\mathrm{A}_2$，$\mathrm{A}_3$は$xy$平面上の点で同一直線上にはないとする．
$3$つの一次式
\begin{align*}
  f_1(x,y)=a_1x+b_1y+c_1, f_2(x,y)=a_2x+b_2y+c_2, f_3(x,y)=a_3x+b_3y+c_3
\end{align*}
は方程式$f_1(x,y)=0$，$f_2(x,y)=0$，$f_3(x,y)=0$によりそれぞれ直線$\mathrm{A_2A_3}$，$\mathrm{A_3A_1}$，$\mathrm{A_1A_2}$を表すとする．
このとき実数$u$，$v$をうまくとると方程式
\begin{align*}
  uf_1(x,y)f_2(x,y)+vf_2(x,y)f_3(x,y)+f_3(x,y)f_1(x,y)=0
\end{align*}
が$3$点$\mathrm{A}_1$，$\mathrm{A}_2$，$\mathrm{A}_3$を通る円を表すようにできることを示せ．

\end{document}

\end{oframed}
\setlength{\columnseprule}{0.4pt}
\begin{multicols}{2}
  {\bf [解答]}

  \begin{figure}[H]
    \centering
    \begin{tikzpicture}[scale=1.25,every node/.style={font=\small}]
      \coordinate (A1) at (-1.2,0);
      \coordinate (A2) at (1.2,0);
      \coordinate (A3) at (0,1.5);
      \coordinate (O) at (0,0.27);
      \draw[very thick,red!85!black] (O) circle (1.23);
      \draw[thick,blue!85!black] ($(A2)!-0.23!(A3)$) -- ($(A2)!1.23!(A3)$);
      \draw[thick,teal!85!black] ($(A3)!-0.23!(A1)$) -- ($(A3)!1.23!(A1)$);
      \draw[thick,orange!90!black] (-1.65,0) -- (1.65,0);
      \node[blue!85!black,right] at (0.75,1.22) {$f_1=0$};
      \node[teal!85!black,left] at (-0.75,1.22) {$f_2=0$};
      \node[orange!90!black,below] at (0,-0.02) {$f_3=0$};
      \node[red!85!black,right] at (1.05,1.55) {$F=0$};
      \foreach \p/\pos in {A1/below left,A2/below right,A3/above}{
        \node[circle,fill=blue!85!black,inner sep=1.2pt] at (\p) {};
        \node[\pos=2pt] at (\p) {$\p$};
      }
    \end{tikzpicture}
    \caption{3直線の交点$A_1,A_2,A_3$を通る円$F(x,y)=0$}
    \label{1998-3:fig:1}
  \end{figure}

  $A_k$ の座標を $A_k(x_k, y_k)$ ($k=1,2,3$) とおく．
  題意の方程式の左辺を
  \begin{align}
    F(x,y)
    = & u f_1(x,y)f_2(x,y) + v f_2(x,y)f_3(x,y)
    + f_3(x,y)f_1(x,y)
  \end{align}
  とおく．

  題意より，$f_1(x,y)=0$ 上に $A_2, A_3$, $f_2(x,y)=0$ 上に $A_1, A_3$, $f_3(x,y)=0$ 上に $A_1, A_2$ があることから
  $F(x,y)=0$ は $3$ 点 $A_k$ ($k=1,2,3$) を通る．

  次に，$F(x,y)$ は高々$2$次で，これが円を表すのは，
  「$x^2, y^2$ の係数が等しく，$xy$ の係数が $0$」
  の時である．したがって$F(x,y)$を展開して$x^2,y^2,xy$の係数比較をして
  \begin{align}
    \begin{dcases}
      (a_1 a_2 - b_1 b_2) u + (a_3 a_2 - b_3 b_2) v + a_3 a_1 - b_3 b_1 = 0 \\
      (a_1 b_2 + a_2 b_1) u + (a_2 b_3 + a_3 b_2) v + a_3 b_1 + a_1 b_3 = 0
    \end{dcases} \label{eq:1}
  \end{align}
  をみたす $(u,v)$ の存在が必要である．
  これらは $u,v$の連立一次方程式だから，解を持つ条件は
  \begin{align}
    \Delta = (a_1 a_2 - b_1 b_2)(a_2 b_3 + a_3 b_2)-(a_3 a_2 - b_3 b_2)(a_1 b_2 + a_2 b_1)
  \end{align}
  が非零であることである．
  整理すると
  \begin{align}
    \Delta
    & = (a_2^2+b_2^2)(a_1 b_3 - a_3 b_1)
  \end{align}
  であり，$a_2, b_2$ は同時には$0$でないから
  \begin{align}
    a_2^2+b_2^2 \neq 0
  \end{align}
  である．また，$f_1(x,y)=0$と$f_3(x,y)=0$が平行ではない直線を表すから
  \begin{align}
    a_1b_3-a_3b_1 \neq 0
  \end{align}
  以上より，
  \begin{align}
    \Delta \neq 0
  \end{align}
  であり，\cref{eq:1}を満たす $(u,v)$ がただ $1$ つ定まる．
  そこでこれを $(u_0, v_0)$ とおく．

  あとはこの値に対して$F(x,y)$が二次式になること，すなわち一次以下でないことを示す．
  $F(x,y)$ の $2$ 次係数が全て $0$ で，$F(x,y)=0$ が $1$ 次方程式だとすると，
  $A_k(k=1,2,3)$ が同一直線上にあることになり矛盾．

  また，直線 $A_3 A_1$ 上の $A_1, A_3$ と異なる点 $(x_0, y_0)$ に対し
  \begin{align}
    F(x_0, y_0) = f_3(x_0, y_0) f_1(x_0, y_0) \neq 0
  \end{align}
  となり，$F(x,y)$ は恒等的には $0$ でない．

  以上から$F(x,y)$は確かに二次式である．

  以上から，$(u,v) = (u_0, v_0)$ とすると $F(x,y)=0$ は円を表す．$\cdots$(答)

  \vspace{10pt}
  {\bf [解説]}

  \newpage
\end{multicols}
\end{document}
